\chapter{Theory}
\label{sec:Theo}
This bachelor's thesis investigates how THz radiation is generated using organic crystals.
This chapter explains the physical foundations underlying this generation.
To begin with, \autoref{sec:NLO}
explains how nonlinear processes are used to generate THz radiation.
In doing so, the susceptibility is not derived completely from quantum mechanics.
After that, \autoref{sec:OR} explains how the THz radiation is generated in the crystal.
In addition, \autoref{sec:PA} explains in more detail why the organic crystals might allow a
greater crystal thickness for generating the THz radiation.
The crystals that were used to generate THz radiation in this bachelor's thesis
are examined in more detail in \autoref{sec:ME}.
Finally, \autoref{sec:EOS} explains the theory behind the measurement method of the setup.


\section{Nonlinear Optics}
\label{sec:NLO}
In optics, most phenomena are linear in nature, meaning that, for example, an electric field $E$
incident on an object is followed by a linear, i.e. equally strong, response, for example of the
polarization density $P$.
However, this linear property is lost once the electric field reaches a strength of
$~3\cdot10^{10}\,\unit{V/m}$.
Such a field strength can only be reached using lasers, which are indeed used here for this
thesis \cite{Butcher_Cotter_1990}.
The nonlinearity becomes apparent when a medium is approximated semiclassically and the linear
solution is shown first.
The medium is approximated as a positive nucleus with an electron attached to it.
An electric field is then applied to this system.
Here, the positive nucleus can be assumed to be spatially fixed, since the electrons are
considerably lighter.
This system can thus be approximated, to first order, as a driven harmonic oscillator.
This yields the equation of motion
\begin{equation}
    \label{eqn:bew}
    m\left(\frac{\text d^2x}{\text d t^2}+2\Gamma \frac{\text d x}{d t}+\Omega^2\cdot x-(\xi^{(2)}x^2+\xi^{(3)}x^3+\dots)\right)= -e E(t)\,.
\end{equation}
Here, $m$ is the mass of the electron, $x$ the displacement from the equilibrium position,
$\Gamma$ the damping constant, $\Omega$ the natural frequency of the system, and $\xi$ the
anharmonic force terms.
In the linear case, only the harmonic solutions are used, and the terms with $\xi^{(2)}$ and
higher are not taken into account.
More details on this can be found in \cite{Butcher_Cotter_1990}.
This is precisely where the nonlinear approach differs, since there the anharmonic terms cannot
be neglected.
Figuratively speaking, in this case the spring connecting the positive nucleus and the electron is
stretched so far apart that the spring's response is no longer linear and edge effects also
become important.
Of particular interest is the polarization $P$, which describes the response of the medium to an
incoming electric field.
Mathematically, this is described as $P(t)= -N\cdot e\cdot x(t)$, with $N$ as the particle
density, $e$ as the charge, and $x(t)$ as the displacement of the electron.
Using the solution of \autoref{eqn:bew} and the definition of the polarization, this can thus be
\begin{equation}
    P =  \epsilon_0 \chi^{(1)}E  \ .
\end{equation}
determined.
In the non-harmonic, i.e. nonlinear, case, however, this cannot be determined exactly, and a
Taylor series expansion is applied.
As a result
\begin{equation}
    \label{eqn:NLP}
    P = \epsilon_0 \left(\chi^{(1)}E+\chi^{(2)}E^2+\chi^{(3)}E^3+\dots\right)
\end{equation}
follows.
Here, $\chi^{(1)}$ describes the linear susceptibility, and the subsequent $\chi^{(2)}$ and
$\chi^{(3)}$ describe the nonlinear susceptibilities of the medium \cite{Butcher_Cotter_1990}.
Here, $\chi^{(2)}$ is a 3rd-rank tensor and additionally has the dependencies $i$, which describes
the polarization, and $j,k$, which are components of the field.
The independence of these components can be reduced through crystal and permutation symmetry.
For the process of optical rectification, the $\chi^{(2)}$ term is decisive.
It is crucial here, however, that this term vanishes in certain cases, for example for a medium
that is inversion-symmetric.
This symmetry property is also referred to as centrosymmetry.
This can be shown using the inversion operation as a symmetry transformation, cf.
\cite[Ch. 5.3.2]{Butcher_Cotter_1990}.
As a result, only 21 of the 32 crystal classes are viable candidates for the generation of THz
radiation.

Furthermore, it is advantageous to consider the polarization in the frequency domain.
For this, it is assumed that the electric field $E$ is composed of several frequency components,
i.e. is defined by
\begin{equation}
    \label{eqn:frq}
    E(t)=E_1e^{i\omega_1 t}+E_2e^{i\omega_2 t} + \dots \ .
\end{equation}
If \autoref{eqn:frq} is then substituted into \autoref{eqn:NLP}, squaring produces further
frequencies such as the sum frequency $\omega_1+\omega2$, the difference frequency
$\omega_1-\omega_2$, and the frequency doublings $2\omega_1,2\omega_2$.
This is called three-wave mixing, since the two input fields and the one output field are
considered \cite{Boyd2023}.
The problem that now arises, however, is that both beams are generated by the same laser and
therefore do not have different frequency components from which to determine their difference.
Here, however, it is crucial that the laser is an ultrashort pulse laser in the fs range, and
therefore, due to the time-bandwidth uncertainty $\Delta t\cdot \Delta \nu \ge \text{const.}$, has
a large frequency bandwidth.
This ultrashort pulse thus has a large bandwidth of frequencies that can mix with one another in
the three-wave mixing process.
The following section shows how THz radiation can be generated via optical rectification, using
difference-frequency generation from a broadband pulse.

% -----------------------978-3-030-73893-8 in Downloads, read TDS chapter 37
\section{Optical Rectification}
\label{sec:OR}


\section{Phase Matching and Coherence Length}
\label{sec:PA}

\section{Material Properties}
\label{sec:ME}

\section{Electro-Optic Sampling}
\label{sec:EOS}
